\chapter{Introduction}
\label{ch:intro}

\section{Clinical Background: Oral Biofilm and Peri-implantitis}
\label{sec:intro_clinical}

Dental implants have become the standard of care for restoring missing teeth, with an estimated 15 million implants placed annually worldwide~\cite{Kolenbrander2010OralMultispecies}.
However, peri-implantitis --- a progressive, biofilm-driven inflammatory condition affecting the peri-implant bone --- affects 20--35\% of implant patients and remains the leading cause of implant failure~\cite{Marsh2003Dental, Berglundh2018Peri, Schwarz2022PeriImplantitis, Sedghi2021PeriImplant}.
The pathogenesis is intimately linked to the multi-species microbial communities that colonise the implant surface: dysbiotic shifts in the oral microbiome, driven by pathobionts such as \textit{Porphyromonas gingivalis} (\textit{Pg})~\cite{Mysak2014Pg}, trigger host immune responses that ultimately lead to bone resorption~\cite{Hajishengallis2012Dysbiosis, Lamont2018Periodontal}.

The oral microbiome is among the most complex in the human body, comprising more than 700 taxa that interact through intricate webs of mutualism, competition, and metabolic cross-feeding~\cite{Kolenbrander2010OralMultispecies, Dewhirst2010HOMD}.
A defining feature of this system is its bistability: in a healthy host, the community is maintained in a commensal state dominated by early colonisers such as \textit{Streptococcus oralis} (\textit{So}) and \textit{Actinomyces naeslundii} (\textit{An}); perturbation can trigger a transition to a dysbiotic state in which late-colonising pathogens --- most notably \textit{Pg}, a low-abundance keystone pathogen~\cite{Hajishengallis2014Polymicrobial} --- establish a positive feedback loop with the immune system~\cite{Marsh2003Dental}.

Two experimental data sources anchor the present work.
First, Heine et al.~\cite{Heine2025PeriImplant} characterised the temporal dynamics of five representative oral taxa (\textit{So}, \textit{An}, \textit{Veillonella} spp.~(\textit{Vei}), \textit{Fusobacterium nucleatum}~(\textit{Fn}), and \textit{Pg}) across four in-vitro conditions defined by community state (commensal vs.\ dysbiotic) and cultivation mode (static vs.\ HOBIC reactor).
The two community states employed distinct \textit{Pg} strains: the commensal consortium used the avirulent strain \textit{Pg} ATCC 20709, whereas the dysbiotic consortium employed the virulent strain \textit{Pg} W83, which harbours a full complement of gingipain and fimbrial virulence determinants.
Second, Dieckow et al.~\cite{dieckow2024} provided a longitudinal in-vivo 16S rRNA amplicon dataset from 10 patients recovering from caries treatment, spanning 10 microbial guilds and three time points over three weeks.
Together, these complementary datasets --- one controlled in-vitro, one clinical in-vivo --- define the scope and the validation strategy of this thesis.


\section{Mathematical Models of Microbial Community Dynamics}
\label{sec:intro_models}

\subsection{Generalised Lotka-Volterra Models}

The generalised Lotka-Volterra (gLV) system, with roots in classical population
ecology~\cite{Lotka1925Elements, May1972Stability}, has long been the workhorse for
modelling microbial community dynamics~\cite{Gonze2018gLV, Stein2013gLV, Venturelli2018Deciphering, Marsland2019MicrobialCRM}:
\begin{equation}
  \dot{\phi}_i = \phi_i \Bigl( r_i + \sum_j A_{ij} \phi_j \Bigr), \quad i = 1,\dots,n,
  \label{eq:glv}
\end{equation}
where $r_i$ is the intrinsic growth rate and $A_{ij}$ the interaction coefficient between species $j$ and species $i$.
Despite its tractability, gLV models suffer from two well-known limitations in the oral context: (i) the interaction matrix $A$ has $n^2$ free parameters, which are severely underdetermined from typical clinical time-course data~\cite{Stein2013gLV, Cao2017Identifiability}; and (ii) the model carries no thermodynamic guarantee of physical consistency --- in particular, there is no constraint preventing non-symmetric interactions that violate energetic reciprocity~\cite{Coyte2015Ecology}. A further subtlety is that relative-abundance (16S) data are inherently \emph{compositional}, so interactions inferred without accounting for the unit-sum constraint can be spurious~\cite{Aitchison1982Compositional, Aitchison1986Compositional, PawlowskyGlahn2015Compositional, Gloor2017Microbiome}.

\subsection{Hamilton Replicator and Extended Hamilton Principle}

A thermodynamically consistent alternative is provided by the Extended Hamilton Principle~\cite{JunkerBalzani2021ExtendedHamilton}, which derives the biofilm evolution equations from a variational free energy functional.
In the homogeneous (0D) limit, this yields the Hamilton replicator equation:
\begin{equation}
  \dot{\phi}_i = \phi_i \Bigl[ (A\bar{\bphi})_i - \bar{\bphi}^\top A \bar{\bphi} \Bigr],
  \label{eq:hamilton_rep}
\end{equation}
where $\bar{\phi}_i = \phi_i \psi_i$ is the living volume fraction, $\psi_i$ is the viability (fraction of living cells within species $i$), and $A$ is \emph{constrained to be symmetric}: $A_{ij} = A_{ji}$.
Symmetry arises automatically from the quadratic form of the free energy and reduces the parameter count from $n^2$ to $n(n+1)/2$; for $n=5$ this gives 15 independent entries.
Biologically, it is motivated by metabolite-mediated interactions where both partners share the same diffusible pool (Section~\ref{sec:theory_replicator}); asymmetric effects such as directional toxin secretion are suppressed and constitute the domain where the unconstrained gLV has a slight predictive advantage.
The mean-fitness subtraction in~\eqref{eq:hamilton_rep} enforces the compositional constraint $\sum_i \phi_i = 1$ exactly, avoiding the spurious volume drift that plagues unconstrained gLV formulations.

\subsection{Parameter Identification: The Inverse Problem}

Estimating the 15-dimensional interaction matrix $A$ from time-course data constitutes a challenging inverse problem.
With $N_t = 5$ usable time points (Day 1 is consumed as an initial condition) and $n = 5$ species, one has at most 25 scalar observations for 15 unknowns; the system is poorly identified without additional constraints.
The posterior landscape is multimodal --- particularly in the dysbiotic condition where the pathogen surge at late time creates near-degenerate solutions --- rendering standard gradient-based optimisation unreliable.

Bayesian inference provides a principled framework for quantifying this uncertainty~\cite{Gelman2013BDA, RobertCasella2004MonteCarlo}; it has recently been applied to biofilm constitutive parameters under hybrid uncertainty~\cite{Fritsch2025BayesianMicrofilms}.
Transitional Markov Chain Monte Carlo (TMCMC)~\cite{ChingChen2007TMCMC, Betz2016TMCMC}, a sequential Monte Carlo sampler~\cite{Chopin2002SMC, DelMoral2006SMCSamplers, Doucet2001SMC}, addresses multimodality through a sequence of annealed intermediate distributions $p_j(\btheta) \propto \mathcal{L}(\mathbf{y}|\btheta)^{\beta_j} p(\btheta)$, where $0 = \beta_0 < \beta_1 < \cdots < \beta_J = 1$, gradually morphing the prior into the posterior. Unlike gradient-based samplers such as HMC/NUTS~\cite{Hoffman2014NUTS, Betancourt2017HMC}, its derivative-free resample--mutate cycle parallelises trivially across particles.
The main computational bottleneck is the repeated forward-model evaluation required at each TMCMC stage; with $N = 2{,}000$ particles and $J \approx 6$ stages, roughly $10^4$--$10^5$ ODE integrations are needed per inference run.

\subsection{Metabolic Priors for Microbiome Inference}

An alternative strategy for reducing ill-posedness is to constrain the \emph{sign} of each interaction coefficient using metabolic network models~\cite{OrthThiele2010WhatIsFBA, Diener2020MICOM, Heinken2023AGORA2}.
If metabolic reconstruction predicts that species $i$ secretes a metabolite that promotes the growth of species $j$ (cross-feeding), then $A_{ij} > 0$ is expected; conversely, competitive exclusion implies $A_{ij} < 0$.
Systematic databases such as AGORA2~\cite{Heinken2023AGORA2} provide genome-scale metabolic reconstructions for thousands of gut and oral taxa, enabling parsimonious FBA (pFBA)~\cite{Lewis2012pFBA} and community FBA (via MICOM~\cite{Diener2020MICOM}) to generate quantitative exchange fluxes from which sign predictions can be derived.
In the second contribution of this thesis (Chapter~\ref{ch:dieckow}), a three-layer sign prior is constructed by combining eHOMD experimental data, AGORA2 pFBA, and MICOM community FBA, and applied to the in-vivo Dieckow dataset.

\subsection{GPU Acceleration in Scientific Computing}

The computational demand of Bayesian inference for ODE-based biofilm models has historically limited its application to low-dimensional problems.
Recent advances in differentiable programming frameworks --- in particular JAX~\cite{jax2018github} --- enable the compilation of the entire forward model to GPU-executable XLA code, with automatic batching of the particle ensemble via \texttt{vmap}.
This approach achieves a ${\sim}200\times$ speedup over the CPU baseline (Section~\ref{sec:heine_gpu}), making TMCMC inference on the 15-dimensional interaction matrix practically feasible.

\subsection{Surrogate Models and Neural Operators}

For further acceleration, neural operator surrogates such as DeepONet~\cite{Lu2021DeepONet}, grounded in the universal-approximation theorem for operators~\cite{Chen1995Universal} and related to Fourier neural operators~\cite{Li2021FNO} and the broader physics-informed and neural-ODE programmes~\cite{Karniadakis2021PINN, Chen2018NeuralODE, Rackauckas2020DiffEq}, can replace the ODE integration in the likelihood evaluation.
DeepONet learns the operator mapping from initial conditions and parameters to the solution trajectory, enabling likelihood evaluations at neural-network inference speed (${\sim}10^3\times$ faster than ODE integration).
Such surrogates are a natural extension of the GPU-accelerated direct-ODE approach taken here and are noted as future work; the present thesis evaluates the forward model exactly.


\section{Objectives and Thesis Structure}
\label{sec:intro_structure}

This thesis makes three main contributions, all grounded in the Extended Hamilton Principle and validated against independent experimental data:

\begin{enumerate}
  \item \textbf{GPU-accelerated Bayesian inference for Heine in-vitro data
        (Chapter~\ref{ch:heine}):}
    A TMCMC framework for estimating the 15-dimensional symmetric interaction matrix
    from Heine et al.'s five-species in-vitro time-course data, achieving ${\sim}200\times$
    speedup via JAX GPU parallelisation. Two-phase estimation (fix-$\psi$ then release-$\psi$)
    resolves posterior multimodality, and UMAP visualisation confirms that all four
    experimental conditions form well-separated clusters in the 15-dimensional
    posterior space.

  \item \textbf{Metabolic-prior replicator dynamics for Dieckow in-vivo data
        (Chapter~\ref{ch:dieckow}):}
    A three-layer metabolic sign prior (eHOMD experimental data, AGORA2 pFBA, MICOM
    community FBA) applied to the Hamilton replicator for in-vivo oral 16S rRNA
    amplicon data from 10 patients. Leave-one-patient-out cross-validation demonstrates
    genuine out-of-sample generalisation (LOO-RMSE $= 0.0490$, Bray-Curtis $= 0.147$,
    47.7\% below the persistence null model), and external validation on Joshi et al.'s
    127 peri-implant samples yields Spearman $\rho = 0.46$ ($p < 0.0001$).

  \item \textbf{Spatial extension, 3-D imaging, and mechanical realisation
        (Chapter~\ref{ch:integration}):}
    A reaction--diffusion extension of the Hamilton model to resolve depth
    stratification, with verified transport numerics (second-order diffusion
    stencil; conservative finite-volume vs.\ non-conservative diagonal scheme), and
    a per-voxel 3-D analysis of $84$ FISH fields of view. The imaging shows that
    dysbiosis is a laterally homogeneous, vertically stratified mat with
    \textit{Pg} sequestered at depth beneath \textit{Fn}
    ($p = 2.2\times10^{-4}$ for the lateral-homogenisation contrast), delimiting what
    a continuum model can and cannot reproduce. A finite-element realisation then
    drives a continuum-mechanics model with this depth composition, linking the
    dysbiotic structure to a higher implant-interface residual stress and a
    $3.3\times$ larger biofilm detachment than the commensal control.
\end{enumerate}

These contributions are deliberately shaped to the implant-research setting that motivates
them. The third in particular closes the loop between the two arms that peri-implant studies
usually pursue separately --- the microbial community colonising the implant surface and the
mechanical fate of the implant itself --- by carrying a 16S-derived compositional dysbiosis
through a finite-element model to an implant-interface mechanical risk; placing the work at
the intersection of peri-implant microbiology and implant biomechanics that defines the field.

The three contributions are unified by the Extended Hamilton Principle
(Chapter~\ref{ch:theory}), which provides the thermodynamic foundation for the
symmetric interaction matrix~$\bA$ and the compositional constraint that
distinguishes Hamilton replicator dynamics from gLV. Across all three, the recurring
finding is that interaction \emph{sign} --- recovered emergently from rich in-vitro
data, imposed as a metabolic prior on sparse in-vivo data, and carried into the
spatial reaction term --- is the robust, transferable structure. Chapter~\ref{ch:integration}
draws out the complementarity of the in-vitro and in-vivo frameworks, sketches a
unified Bayesian--metabolic inference scheme, and reports the spatial results before
closing on clinical applications.
